% ============================================================
% Section 6: Conclusion
% ============================================================

\section{Conclusion}
\label{sec:conclusion}

% \section{おわりに}
% \label{sec:conclusion}
%
% 本章では，本研究の主要な貢献を総括し，限界と今後の研究方向を述べる．

% ============================================================
\subsection{Summary of Contributions}
\label{sec:summary}
% ============================================================

% \subsection{貢献のまとめ}
% \label{sec:summary}

This paper proposed and evaluated \emph{Preemptive Owner-Priority Scheduling}
for the campus GPU sharing platform Air Computing Pool (ACP).
The core problem motivating this work is that existing sharing policies
(FCFS and Owner Priority) fail to protect high-performance GPU owners
at high system load, causing rational owners to withdraw and triggering
a system-wide participation cascade collapse.
We showed that preemption, reinterpreted not as a throughput tool
but as an \emph{ownership insurance mechanism}, eliminates this failure mode.

% 本論文は，学内GPU共有基盤Air Computing Pool（ACP）に対して
% \emph{プリエンプティブ所有者優先スケジューリング}を提案・評価した．
% 本研究の動機となる核心的問題は，既存共有方式（FCFSおよび所有者優先）が
% 高負荷時に高性能GPU所有者を保護できず，合理的な所有者が離脱を選択して
% システム全体の参加カスケード崩壊を引き起こすことである．
% プリエンプションをスループットツールとしてではなく
% \emph{所有権保険メカニズム}として再解釈することが，この失敗モードを解消することを示した．

The three main contributions of this work are:

% 本研究の3つの主要貢献は以下のとおりである：

\textbf{(1) Problem reformulation.}
We identified that the fundamental challenge in campus GPU sharing is not
throughput maximization, but maintaining the \emph{participation incentive}
of high-performance GPU owners.
Without owner participation, the shared pool degrades and the entire
ecosystem collapses, regardless of how efficiently the scheduler operates.

% \textbf{(1) 問題の再定義．}
% 学内GPU共有における根本的な課題は，スループット最大化ではなく，
% 高性能GPU所有者の\emph{参加インセンティブ}を維持することであることを特定した．
% 所有者が参加しなければ，スケジューラの効率性に関わらず，
% 共有プールは劣化してエコシステム全体が崩壊する．

\textbf{(2) Preemption as ownership guarantee.}
We showed that immediate preemption upon owner arrival guarantees
PR $\leq 1$ at all load levels by construction:
the owner's TAT under sharing is always at most the standalone TAT.
This converts ACP participation from a load-dependent gamble
into a dominant strategy for rational owners.

% \textbf{(2) 所有権保証としてのプリエンプション．}
% 所有者到着時の即時プリエンプションが，設計上すべての負荷水準でPR $\leq 1$を保証することを示した：
% 共有下での所有者TATは常に単独利用時のTAT以下である．
% これにより，ACP参加を負荷依存の賭博から合理的所有者にとっての支配戦略へと転換する．

\textbf{(3) Empirical validation across diverse workloads.}
Through 100-trial discrete-event simulations covering uniform and heterogeneous
workload distributions (Low-Heavy, High-Heavy, Random training ratios),
we demonstrated that Preemptive is the only policy that:
(a)~maintains PR $\leq 1$ in three out of four scenarios across all load levels,
(b)~achieves a self-sustaining participation equilibrium without external
incentives from any initial participation state, and
(c)~provides substantially lower TAT for low-tier users
compared to No Sharing.

% \textbf{(3) 多様なワークロードにわたる経験的検証．}
% 均一および不均一なワークロード分布（Low-Heavy・High-Heavy・Randomのtraining ratio）を
% カバーする100試行の離散イベントシミュレーションを通じて，
% プリエンプティブ方式が唯一以下を満たす方式であることを実証した：
% (a)~4シナリオ中3つで全負荷水準においてPR $\leq 1$を維持する，
% (b)~外部インセンティブなしに，任意の初期参加状態から自律的な参加平衡を達成する，
% (c)~No Sharingと比較して低性能層ユーザーに大幅に低いTATを提供する．

% ============================================================
\subsection{Limitations}
\label{sec:limitations}
% ============================================================

% \subsection{限界}
% \label{sec:limitations}

Several limitations of the current study should be noted.

% 現在の研究のいくつかの限界を指摘する．

\textbf{Simulation fidelity.}
The simulation models task arrival as a Poisson process and task sizes
as log-normal distributions fitted to the Alibaba cluster trace.
Real workloads may exhibit burstiness, long-range dependence, or
correlated arrivals that deviate from these assumptions.
The checkpoint overhead factor $\alpha = 0.2$ is a conservative estimate;
actual values depend on GPU model, task memory footprint, and network bandwidth.

% \textbf{シミュレーションの忠実度．}
% シミュレーションはタスク到着をポアソン過程，タスクサイズをAlibaba cluster traceに
% フィットした対数正規分布としてモデル化している．
% 実際のワークロードは，これらの仮定から外れるバースト性・長距離依存性・
% 相関到着を示すことがある．
% チェックポイントオーバーヘッド係数$\alpha = 0.2$は保守的な推定値であり，
% 実際の値はGPUモデル・タスクメモリフットプリント・ネットワーク帯域幅に依存する．

\textbf{Scale.}
The evaluation is limited to 18 users and 9 GPU tiers.
Behavior in larger deployments (hundreds of users, heterogeneous hardware
generations) remains to be characterized.

% \textbf{スケール．}
% 評価は18ユーザー・9GPUティアに限定される．
% より大規模な展開（数百ユーザー，異種ハードウェア世代）での挙動は未評価である．

\textbf{Low-Heavy exception.}
In the Low-Heavy scenario at $\rho_\mathrm{sys} \geq 0.9$,
checkpoint overhead from large guest training jobs causes the Preemptive
PR to slightly exceed 1.0 for inference-dominant owners.
Although this violation is marginal and confined to near-saturation loads
uncommon in practice, it shows that the PR $\leq 1$ guarantee is not
unconditional.

% \textbf{Low-Heavyの例外．}
% Low-Heavyシナリオの$\rho_\mathrm{sys} \geq 0.9$において，
% 大規模ゲスト学習タスクのチェックポイントオーバーヘッドにより，
% プリエンプティブのPRが推論主体の所有者に対してわずかに1.0を超える．
% この違反は軽微で実際には稀な飽和負荷に限定されるが，
% PR $\leq 1$保証が無条件ではないことを示している．

\textbf{Single incentive mechanism.}
The model considers only ownership (TAT-based utility) as the incentive
for participation.
Real environments may involve additional incentives such as monetary rewards,
reputation, or administrative mandates that could interact with
the ownership guarantee in complex ways.

% \textbf{単一インセンティブメカニズム．}
% モデルは参加インセンティブとして所有権（TATベースの効用）のみを考慮している．
% 実際の環境では，金銭的報酬・声望・管理的指令などの追加インセンティブが
% 所有権保証と複雑な形で相互作用する可能性がある．

% ============================================================
\subsection{Future Work}
\label{sec:future_work}
% ============================================================

% \subsection{今後の研究方向}
% \label{sec:future_work}

Several directions for future work are identified.

% いくつかの今後の研究方向を特定する．

\textbf{Implementation and real-world evaluation.}
The next step is to implement Preemptive Owner-Priority Scheduling
in a production GPU cluster scheduler (e.g., Kubernetes with GPU plugins,
or SLURM with GRES).
Real-world measurements of preemption latency, checkpoint size, and
migration cost would allow calibrating $\alpha$ per GPU model and
workload type, and validating the simulation results at scale.

% \textbf{実装と実環境評価．}
% 次のステップは，プリエンプティブ所有者優先スケジューリングを
% 本番GPUクラスタスケジューラ（GPUプラグイン付きKubernetes，GRES付きSLURM等）に実装することである．
% プリエンプションレイテンシ・チェックポイントサイズ・マイグレーションコストの実測により，
% GPUモデル・ワークロード種別ごとに$\alpha$を較正し，
% シミュレーション結果のスケールでの検証が可能となる．

\textbf{Reducing preemption overhead for inference owners.}
To address the Low-Heavy limitation, future work could explore
\emph{kill-and-restart} preemption for short inference guest tasks
(where restarting from scratch is cheaper than checkpointing) and
\emph{migration-free preemption} strategies that allow the owner's
inference job to start on a temporarily idle GPU while the
checkpointed guest task is re-queued.

% \textbf{推論所有者向けプリエンプションオーバーヘッドの削減．}
% Low-Heavyの限界に対処するため，短い推論ゲストタスクに対する
% \emph{kill-and-restart}プリエンプション（チェックポイントよりも最初から再起動が安価な場合）と，
% チェックポイントされたゲストタスクが再キューイングされる間に
% 所有者の推論ジョブを一時的に遊休状態のGPUで開始する
% \emph{migration-free preemption}戦略の探索が考えられる．

\textbf{Scalability.}
Future work should evaluate the proposed method in larger environments
(hundreds of users, multiple GPU generations) and design distributed
scheduling algorithms that maintain the per-GPU ownership rule
without requiring centralized coordination.

% \textbf{スケーラビリティ．}
% 将来の研究では，より大規模な環境（数百ユーザー，複数GPUバージョン）での
% 提案方式の評価と，集中型コーディネーションなしにGPU単位の所有権ルールを維持する
% 分散スケジューリングアルゴリズムの設計が必要である．

\textbf{Combined incentive mechanisms.}
Integrating ownership-based protection with complementary incentives
(e.g., usage credits proportional to GPU donation time, or
priority boosts for frequent contributors) could further strengthen
the participation equilibrium and attract users who are indifferent
to ownership alone.

% \textbf{複合インセンティブメカニズム．}
% 所有権ベースの保護と補完的インセンティブ
% （例：GPU寄付時間に比例した利用クレジット，頻繁な貢献者への優先度ブースト）の統合により，
% 参加平衡をさらに強化し，所有権だけでは無関心なユーザーを引き付けることができる．

\textbf{Broader applicability.}
The ownership-guarantee model introduced in this paper is not specific
to GPU sharing.
The same framework---in which a contributed resource carries an
instant-reclaim preemption right---can be applied to other shared campus
resources such as CPU clusters, storage pools, or network bandwidth,
and to peer-to-peer resource sharing in enterprise or cloud environments.

% \textbf{より広い適用可能性．}
% 本論文で導入した所有権保証モデルはGPU共有に限定されない．
% 貢献リソースが即時回収プリエンプション権を持つという同じフレームワークは，
% CPUクラスタ・ストレージプール・ネットワーク帯域幅などの
% 他の学内共有リソースや，企業・クラウド環境のP2Pリソース共有にも適用できる．
